\begin{abstract}
\noindent
I study how the cost of family space shapes fertility and how can housing
policy affect it. Housing costs are at the center of the public debate
on declining fertility: children need space, and young households, who have
little savings, find that space costly to pay for. I build a life-cycle model
in which households choose whether to have children, whether and how much to
rent or own, and how much to save. Buying requires a down payment, and
earnings are risky. I calibrate the model to the United States in 2007 and
treat the fall in births since then as an exogenous decline in the taste for
children. Along the transition to a new steady state with a smaller
population, fewer households make housing cheaper, which feeds back into
fertility. Because a child is a long commitment, earnings risk weighs on
fertility, and lower risk raises births. I then use the model to assess
housing policy. Easier mortgage credit could give young households access to
the larger homes a child requires, but it does little for births. Policies
that change long-term incentives, such as more housing supply and credit for
parents, raise fertility, and their effects build up over a generation.
Housing policy must therefore be judged by its long-run effects.

\end{abstract}
